Given a prime p, let ord_p(a) denote the multiplicative order of a modulo p, i.e.
ord_p(a) is the least positive integer t such that \(a^t \equiv 1 (mod p)\), assuming p \nmid a.
An integer a is a primitive root modulo p if ord_p(a)=p-1.

Literal statement (as written in the problem file):

  (S_lit)  For every prime p, there exists a prime q<p such that ord_p(q)=p-1.

Immediate ambiguity/edge case:
- If p=2, there is no prime q<p, so (S_lit) is false for the literal quantifier "every prime".

Minimal corrected statement (the version that seems intended in the number-theoretic context):

  (S)  For every odd prime p, there exists a prime q<p such that ord_p(q)=p-1.

Everything below separates (S_lit) (which is trivially false) from (S) (which appears open).
